% Testkorpus: mehrseitiges Datenblatt mit Tabellenumbruch, Einheiten,
% Kopf-/Fußzeilen und Seitenzahl „x von y“.
\documentclass[a4paper,10pt]{article}
\usepackage{fontspec}
\usepackage[ngerman]{babel}
\usepackage[margin=2cm]{geometry}
\usepackage{longtable,tabularx,array,booktabs,multirow}
\usepackage[locale=DE]{siunitx}
\usepackage{fancyhdr}
\usepackage{lastpage}
\usepackage{xcolor}
\usepackage{colortbl}
% 120 Tabellenzeilen erzwingen Seitenumbrüche in der longtable.
\newcounter{nr}
\newcommand\Zeile{\stepcounter{nr}\arabic{nr} & Punkt & \SI{230}{\volt} \\}
\newcommand\Zehn{\Zeile\Zeile\Zeile\Zeile\Zeile\Zeile\Zeile\Zeile\Zeile\Zeile}
\pagestyle{fancy}
\fancyhead[L]{Datenblatt}
\fancyfoot[C]{Seite \thepage\ von \pageref{LastPage}}
\begin{document}
\section*{Messwerte}
\begin{tabularx}{\linewidth}{lX S[table-format=3.2]}
\toprule
Größe & Beschreibung & {Wert [\si{\milli\metre}]} \\
\midrule
A & Länge & 12.50 \\
B & Breite & 7.25 \\
\bottomrule
\end{tabularx}

\begin{longtable}{r l r}
\toprule
Nr. & Messpunkt & Spannung \\
\midrule
\endhead
\bottomrule
\endfoot
\Zehn\Zehn\Zehn\Zehn\Zehn\Zehn\Zehn\Zehn\Zehn\Zehn\Zehn\Zehn
\end{longtable}
\end{document}
